\clearpage\thispagestyle{empty}
\pdfbookmark[0]{Quellen zu Gauß an Bode (26. Dezember 1821)}{bodedecember1821-sources}
\begin{center}\textbf{Quellen zu Gauß an Bode (26. Dezember 1821)}\end{center}
\noindent Diese KI-Übertragung ins heutige Deutsch folgt dem Brief vom 26. Dezember 1821 in \href{https://gdz.sub.uni-goettingen.de/id/PPN236007467}{Gauß’ \textit{Werken}, Bd.~VI} (Göttingen, 1874), S.~439--440. Johann Elert Bode druckte ihn unter der Überschrift \textit{Erfindung eines Heliotrops, Beobachtungen des Kometen von 1821 und Elemente seiner Bahn} im \textit{Astronomischen Jahrbuch für das Jahr 1825}, Bd.~50, S.~\href{https://api.digitale-sammlungen.de/iiif/image/v2/bsb10538354_00109/full/1600,/0/default.jpg}{103}, \href{https://api.digitale-sammlungen.de/iiif/image/v2/bsb10538354_00110/full/1600,/0/default.jpg}{104}, \href{https://api.digitale-sammlungen.de/iiif/image/v2/bsb10538354_00111/full/1600,/0/default.jpg}{105} und \href{https://api.digitale-sammlungen.de/iiif/image/v2/bsb10538354_00112/full/1600,/0/default.jpg}{106}. Das \href{https://api.digitale-sammlungen.de/iiif/image/v2/bsb10538354_00003/full/1600,/0/default.jpg}{Titelblatt} des Exemplars der Bayerischen Staatsbibliothek nennt Berlin 1822. Im Impressum stehen Bode als Verleger, Ferd. Dümmler als Kommissionär und C. F. E. Späth als Drucker. Die Quellenzeile der Werkausgabe nennt Berlin 1821 und S.~103--105. Die Vergleichstabelle füllt S.~\href{https://api.digitale-sammlungen.de/iiif/image/v2/bsb10538354_00112/full/1600,/0/default.jpg}{106} des Jahrbuchs.
\par\smallskip\noindent Im Brief stehen beim ersten Zeitwert 42 Sekunden, ebenso im Jahrbuch. Im früheren \hyperlink{native-442-part-lower}{Bericht vom 17. Mai 1821} stehen 32 Sekunden. Beide Werte bleiben an ihrer jeweiligen Stelle erhalten. Das Jahrbuch nennt ebenfalls $9\tfrac{1}{3}$ geografische Meilen. Gauß schreibt die Berechnung der Elemente von Staudt zu und legt dessen eigenhändigen Vergleich bei.
\par\smallskip\noindent Auf S.~\href{https://api.digitale-sammlungen.de/iiif/image/v2/bsb10538354_00110/full/1600,/0/default.jpg}{104} versah Bode den Verweis auf die Mitteilung über den Heliotrop mit einer unterzeichneten Fußnote. In heutiger Schreibung lautet sie „Siehe im letzten Abschnitt \textit{Noch verschiedene astronomische Beobachtungen, Bemerkungen usw.} B.“ Die Werkausgabe verweist an dieser Stelle redaktionell auf Band IV. KI-Übertragung und KI-Satz.
\clearpage
